% =====================================================================
%  Chapter 9 — Step 3
% =====================================================================

\chapter{Step Three: The Sign-Only Weil Argument}
\label{ch:step3}

This is the step where the argument stops resembling its ancestor and starts
doing something new. Everything so far has been preparation; this is where the
Weil positivity is supposed to come out.

The remark's item 5 is the longest sentence in the document, and it is worth
reproducing in full before unpacking it, because the reader needs to see how
compressed it is:

\begin{quote}
\small
``definitively switch to $\Lambda$ aka $\sum_{n \le X}\Lambda(n)\cdot(X/n-1)_+ -
X\log X + (1+\gamma)\cdot X$; Guinand--Weil had the sign, variance test function
trick to $\sum_{h,x_1,\dots,x_k} t(x_1)\cdots t(x_k)\,e_q(h(x_1x_2 + \cdots +
x_{k-1}x_k))$, sum over $h$ first for the left-hand side, with e.g.\ quadric
coefficients too Bertini gives a definite sign from middle cohomology, one off
the one on $M_k$, which is the zero-mode on the right-hand side with $R$
square-root-quality up to an inch or so as Katz would say by iterating Parseval
$k$ times, opposite signs meaning no cancellation of inverse roots, QED;
vectorize to go up dimensions.''
\end{quote}

Four distinct pieces of mathematics are compressed in there. This chapter
takes them one at a time.

\section{Step 3a: switch to the von Mangoldt sum}The first move is to change the object being bounded. Instead of working with
$\zeta$ directly, the argument works with
\begin{equation}
\Psi(X) = \sum_{n \le X}\Lambda(n)\left(\frac{X}{n} - 1\right)_{+} - X\log X + (1+\gamma)X,
\label{eq:psi}
\end{equation}
which is the remark's $\Lambda$-switch and, up to the standard terms, the
explicit formula for $\psi(X)$ with the zeros replaced by their contribution to
the error term.

Here is why this is the right object. Landau's oscillation theorem says
$\psi(x) - x = \Omega_\pm(x^{\Theta - \varepsilon})$, with $\Theta$ the
supremum of the real parts of the zeros. The Riemann hypothesis is the
statement $\Theta = \tfrac12$. So one wants to bound the *oscillation* of
$\psi$ around its main term $x$, and the explicit formula says that oscillation
is exactly the sum over the zeros of $\zeta$ of terms of size $x^{\rho}$, with
$x^{1/2}$ the size of a critical-line zero.

Two features of \eqref{eq:psi} are worth noting.

\begin{description}[style=nextline,leftmargin=0pt,font=\sffamily\bfseries\small,
  itemsep=6pt]
\item[The truncation $(X/n-1)_+$ is the smoothing.] The explicit formula for
$\psi$ is a distribution statement; the smoothed version replaces the
$\delta$-functions at the zeros by a smooth kernel, and the answer is then a
sum over the zeros of a smooth function of the ordinates.
\item[The $(1+\gamma)X$ term is the residue at $s=1$.] It is there because
$\zeta$ has a pole; without it the formula would be off by a linear term.
\end{description}

The switch is not a shortcut. It is the standard analytic form of the problem:
RH is a bound on $\Psi(X)$, and the Weil expression $W$ is the quadratic form
that encodes exactly such a bound.

\section{Step 3b: the test-function trick, and the sum over h first}The remark's phrase ``Guinand--Weil had the sign'' refers to the classical
asymmetry in the moments of $\zeta$ on the critical line, refined by Guinand:
for the right test function, the moments of $\zeta$ on the line are smaller
than those off it, and the difference has a sign. That sign is what a
positivity argument can use.

The "variance test function trick" is the analytic engine. One chooses a test
function $t$ and forms
\begin{equation}
\sum_{h,\,x_1,\dots,x_k} t(x_1)\cdots t(x_k)\,e_q\!\left(h(x_1x_2 + \cdots + x_{k-1}x_k)\right),
\label{eq:trick}
\end{equation}
where $e_q$ is the additive character of $\mathbb{F}_q$, and the sum runs over
the coefficients $h$ of a quadratic form and over the coordinates $x_i$.

The instruction that makes it work is ``sum over $h$ first''. A quadratic form
in $k$ variables over $\mathbb{F}_q$ has an associated character sum, and
summing over the coefficients $h$ first replaces the sum by a count of solutions
of a quadratic equation --- a *point count*. And a point count is arithmetic.
That is the moment the method reconnects to the function-field ancestor: the
quadratic form supplies the geometry, and the count of its solutions supplies
the arithmetic factor.

The remark's parenthetical is decisive here: ``with e.g.\ quadric
coefficients too Bertini gives a definite sign from middle cohomology''. A
quadric in $k$ variables has a very well-understood middle cohomology, and
that cohomology determines the sign of the character sum. So the sign of the
sum over $h$ is a cohomological fact, computable, and --- crucially --- the
\emph{same} over $\mathbb{Z}$ if the construction is automorphic rather than
geometric over $\mathbb{F}_q$.

That is the "natively-automorphic" claim of the remark, stated in the place
where it does the work.

\section{Step 3c: the identity qN = Mk + R}The remark's item 3 gives the algebraic form:

\begin{quote}
\small
``$qN = Mk + R$, using $x_1\cdot x_2 + \cdots + x_{k-1}\cdot x_k$, wanting
`automorphy' of $N$ since $\mu$ is automorphic from $\theta$ on sixfold cover
of $\mathrm{GL}_5$ as a period integral of thetas, quadric has natural
interpretation on corresponding CM e.g.\ $\mathbb{C}/\mathbb{Z}[\omega]$ as
Weil, put Poincar\'e bundle on the edges to get $\theta$s.''
\end{quote}

So the identity is a three-term decomposition
\begin{equation}
qN = Mk + R,
\label{eq:qNMR}
\end{equation}
in which:

\begin{description}[style=nextline,leftmargin=0pt,font=\sffamily\bfseries\small,
  itemsep=6pt]
\item[$q$] is the prime power, the arithmetic factor: the size of the field over
which the character sum is taken.
\item[$N$] is the number being counted.
\item[$M$] is the main term: the expected count, of size $q^{\deg}$, coming
from the trivial solutions.
\item[$R$] is the remainder: the excess over the main term, which is what
carries the information about the zeros.
\end{description}

The coefficient $M_k$ is the count of the trivial solutions of the quadratic
equation in $k$ variables, and it depends only on $k$. The identity is the
statement that the total count splits into a main term that is completely
understood and a remainder that is not.

The method's target is the remainder. And the requirement on $R$ is exactly the
requirement for the Weil bound to be sharp enough: $R$ must have
\emph{square-root quality}. The remark's phrase is ``$R$ square-root-quality up
to an inch or so as Katz would say by iterating Parseval $k$ times'', and
``up to an inch or so'' is the author being careful: the bound is
$R = O(q^{k/2 + \varepsilon})$ up to a bounded loss, not exactly
$O(q^{k/2})$.

\subsection{Where the square root comes from}The reason a remainder of square-root quality is the right target is the
Weil bound on a curve: the excess in the count of points over a curve of genus
$g$ is $O(q^{1/2})$, with $q$ the base field and the number of points
$q + 1 + O(g\sqrt q)$. If the geometry in the sum over $h$ is that of a curve,
then the remainder of a character sum over that curve is of square-root size
by the Weil bound for the relevant $L$-function.

So the passage from the function-field method to the integer method has the
following form, and this is the conceptual heart of the whole claim:

\begin{quote}
The Weil bound gives $O(\sqrt{\cdot})$ for the excess over the main term in a
point count, because the point count is the trace of Frobenius and the Weil
bound is the bound on the eigenvalues of Frobenius. Over the integers there is
no Frobenius and no point count --- but there \emph{is} a Weil expression,
which is the same quadratic form, and its non-negativity is the same
statement. If the remainder $R$ can be shown to have square-root quality by an
automorphic rather than a geometric argument, then the Weil bound has been
reproduced, and with it the hypothesis.
\end{quote}

That is why the argument needs the sign and the size together: the sign gives
the Weil positivity (RH), and the size gives the sharp bound. And it is why
the determinant matters: a determinant is how one writes a quadratic form as a
Gram matrix, and a Gram matrix is how one certifies non-negativity without
computing the sum.

\section{Step 3d: "opposite signs meaning no cancellation of inverse roots"}The remark's phrase ``opposite signs meaning no cancellation of inverse
roots'' is the technical heart of the sign argument, and it is the part that
most rewards a slow reading.

The idea is this. The character sum $\eqref{eq:trick}$ is a sum of terms
$\pm q^{1/2}\cdot(\text{root})$, where the roots are the reciprocals of the
eigenvalues of Frobenius on the cohomology of the quadric. If two terms with
opposite signs had the same root, they would cancel, and the bound would
degrade: the sum would be as large as the number of terms rather than as large
as the square root of the number of terms. The bound holds only if the
cancellation cannot happen --- that is, if the terms with opposite signs have
\emph{different} roots.

To force this, one iterates. The remark says ``by iterating Parseval $k$
times'': one applies the Parseval identity to separate the sum by the sign,
then applies it again to separate each sign class by the root, and continues
for $k$ iterations. After $k$ iterations, the sign and the root are separated,
and one reads off that the roots attached to opposite signs are disjoint.

This is the classical ``orthogonality of the columns of the integral
structure'' that the remark's item 16 names, and it is where the requirement
of an integral structure across Hecke characters earns its keep.

It is also, not coincidentally, where the reference to $\mathrm{GL}_5$ comes
in. The remark's item 16 says the automorphy of $\mu$ comes from a sixfold
cover of $\mathrm{GL}_5$ and that the argument then reduces, via
Cauchy--Schwarz, to a bound on the norm of a suitably projected Poincar\'e
series, using the approximate orthogonality of the columns of the integral
structure.

\begin{quote}
The formalisation carries a statement that is exactly the separation the remark
describes. In
\code{Formalism.SharpRestriction.Divisibility.FunctionModelRallisFromInputs}
the condition \code{NCGaps} reads
{\footnotesize
\begin{verbatim}
forall j, j <= D.R.m -> forall tau, forall x in P.Dpar j
  tau,
  forall y in P.Dpar j tau, x.2 != y.2 ->
    ((D.R.r + D.R.m * P.t : Nat) : Q)^3 < |x.1 - y.1|
\end{verbatim}
}
i.e.\ the separation between distinct roots is required to exceed a
\emph{cube} of the weight-and-level combination. A cube, not a square: the
extra power is the cost of the $k$-fold Parseval iteration that separates
sign from root.
\end{quote}

This is worth flagging as a methodological point. The remark's own summary
says ``opposite signs meaning no cancellation of inverse roots''; the
formalised condition is a quantitative separation with a cube. The two are
consistent, and the formalised one is much more informative. A reader who
wants to evaluate the argument should work from the formalised condition, not
from the sentence.

\section{Step 3e: the GGP periods and the interlacing weights}The remark's item 7 supplies the machinery for the determinant:

\begin{quote}
\small
``Determinants for $W(\cdot)$'s containing the desired periods, indeed GGP
periods, with values 0, 1 meaning interlacing weights, or up to $m$
non-interlacing ones.''
\end{quote}

The \emph{Gross--Prasad periods} are the periods of the rank-difference
Shimura variety attached to a pair of groups; they are the natural
automorphic periods for the unitary groups in play. The remark needs them as
the entries of the matrix whose determinant is the Weil expression.

The interlacing condition is the technical heart. Suppose the Weil expression
is written as a determinant of a matrix whose rows are indexed by the
ordinates of the zeros and whose columns by the periods. A row whose entries
come from periods whose arguments interlace the zero's argument contributes a
sign; a determinant built out of interlacing data is a sum of terms of
alternating sign, and the vanishing of the alternating sum is the positivity.

This is the classical mechanism behind the Weil bound for a curve --- the
number of points of a curve over $\mathbb{F}_q$ is $q + 1 - \sum_\alpha
\alpha$, and the $\alpha$ are the Frobenius eigenvalues, whose arguments
interlace. The claim of the remark is that the same interlacing, with
automorphic periods in place of Frobenius eigenvalues, can be manufactured over
$\mathbb{Z}$.

\begin{seamnote}{What the development says about this}
\code{Development/WeilBound/ThetaLifts/FrameEstimates/} is where the frames
and their estimates live --- \code{FrameBasis.lean},
\code{ExistsColsDetNeZero.lean} (the non-vanishing of the relevant
determinant), \code{BlockForm.lean}, \code{DomWeil.lean} (domination of the
Weil expression), \code{SigRapid.lean} and \code{CuspRapid.lean} (the rapid
decay of the signature and of the cusp form). The final assembly is
\code{sigRapid\_of}, which takes the inputs
\code{ThetaBasic}, \code{UMGLin}, \code{EpsOf}, \code{IUMG}, \code{IInv},
\code{TriplePair} and \code{DomWeil} and returns the rapid signature estimate.

Note the shape: the estimate is a consequence of a \emph{bundle} of
ingredients, each separately proved, and the Weil domination \code{DomWeil} is
itself an input to the final theorem rather than something derived inside it.
That is the same bundling pattern that appears in
Chapter~\ref{ch:inputs}, and it is the pattern that makes the statement review
of that chapter necessary.
\end{seamnote}

\section{Step 3f: vectorising to go up dimensions}The last phrase of the remark's item 5 is ``vectorize to go up dimensions''.

The Weil expression $W$ is a quadratic form in a scalar function of the
ordinates. To handle it in $n$ dimensions, or to handle the higher-rank
groups that the theta lift moves between, one needs the vector-valued
analogue: a quadratic form with values in a vector space of representations,
and a determinant of a matrix of such forms. The remark's author, in the
footnote to item 5, describes the cost of this in his own terms:

\begin{quote}
\small
``A lot is obviously being lossily compressed here, but especially after
gritting with Sussman through Ba\'laban's fantastic induction I had to think
to myself about how hard fighting for some sort of $\mathbb{F}_1$ \'a la e.g.\
Borger would be.''
\end{quote}

The reference is to work on $\mathbb{F}_1$-type foundations for the trace
formula and for ``$\mathbb{F}_1$-like'' induction, and to the difficulty of
making that rigorous. That is a real and known difficulty: the $\mathbb{F}_1$
paradigm is a heuristic, not a theory, and the remark's author is aware that
he is using it heuristically.

\section{What step 3 has and has not shown}Step 3, if it works, establishes the Weil positivity and hence RH. It rests
on:

\begin{enumerate}
\item the Weil criterion's equivalence of RH with a positivity;
\item an automorphic replacement for the character sum over $\mathbb{F}_q$;
\item a determinant whose Gram structure certifies the positivity;
\item a separation of sign from root, obtained by $k$-fold Parseval iteration,
formalised as a cubic gap condition.
\end{enumerate}

Item 2 is the one that the remark calls the ``natively-automorphic
determinant'', and it is the one with no classical precedent.

\begin{keytakeaways}
\item The method switches to the von Mangoldt sum, where RH becomes a bound
on the oscillation of $\psi$ around $x$, and the Weil expression $W$ encodes
exactly that bound.
\item The test-function trick sums over the coefficients of a quadratic form
first, which converts the sum into a point count --- the arithmetic factor that
the integers lack --- and the sign of that count comes from the middle
cohomology of the quadric.
\item The identity $qN = Mk + R$ separates the understood main term from the
informative remainder, and the requirement on $R$ is square-root quality.
\item The positivity is certified by a determinant, because a quadratic form is
non-negative exactly when it has a Gram factorisation, and the Gram matrix is
built out of GGP periods with interlacing weights.
\item The separation of sign from root is what keeps the bound sharp; it is
formalised as a cubic separation condition, and it is the price of the
$k$-fold Parseval iteration.
\end{keytakeaways}
